% ===================================================================
% CORE DOCUMENT UTILITIES
% ===================================================================
\usepackage[margin=0.5in]{geometry}   
\usepackage{graphicx}                 
\usepackage{float}                   
\usepackage{booktabs}                 
\usepackage{array}                    
\usepackage{multicol}                 
\usepackage{parskip}                  
\usepackage{titling}                  
\setlength{\droptitle}{-1cm}          


% --- Blank Page Command ---
\newcommand*{\blankpage}{
  \begingroup
    \setlength{\topskip}{0pt}
    \setlength{\parskip}{0pt}
    \vspace*{\fill}
    \centering
    \nointerlineskip 
    This page would was intentionally left blank. \par
    \vspace{\fill}
  \endgroup
}

% --- TODO ---
\newcommand{\todo}[1]{\textcolor{red}{\textbf{[TODO: #1]}}}

% ===================================================================
% COLORS, TEXT, AND LISTS
% ===================================================================
\usepackage[svgnames]{xcolor}
\usepackage{soul}             
\usepackage[inline]{enumitem}         
\usepackage{gensymb}                  
\usepackage{xspace}                  

% ===================================================================
% MATHEMATICS, LOGIC, AND SEMANTICS
% ===================================================================
\usepackage{proof}             
\usepackage{amsmath, amssymb, mathtools} 
\usepackage{mathrsfs}                   
\usepackage{stmaryrd}                    
\usepackage{mathpartir}                  
\usepackage{semantic}                    
%\usepackage{proofpack}             
\usepackage{empheq}      

% --- Paired delimiters ---
\newcommand{\set}[1]{\left\{#1\right\}}
\newcommand{\abs}[1]{\left\lvert#1\right\rvert}
\newcommand{\norm}[1]{\left\lVert#1\right\rVert}
\newcommand{\paren}[1]{\left( #1 \right)}
\newcommand{\brak}[1]{[#1]}
\newcommand{\floor}[1]{\left\lfloor#1\right\rfloor}
\newcommand{\ceil}[1]{\left\lceil#1\right\rceil}

% --- Logic and Semantics ---
\newcommand{\Iof}[1]{\mathcal{I}(#1)}
\newcommand{\Ivof}[1]{\mathcal{I}_{\varphi}(#1)}
\newcommand{\sat}{\vDash}
\newcommand{\unsat}{\nvDash}
\newcommand{\satphi}{\vDash^{\varphi}}
\newcommand{\unsatphi}{\nvDash^{\varphi}}
\newcommand{\der}{\vdash}
\newcommand{\derC}{\vdash_{\cC}}
\newcommand{\notder}{\nvdash}


% --- Calligraphic uppercase letter ---
\newcommand{\cA}{\mathcal{A}}
\newcommand{\cB}{\mathcal{B}}
\newcommand{\cC}{\mathcal{C}}
\newcommand{\cD}{\mathcal{D}}
\newcommand{\cE}{\mathcal{E}}
\newcommand{\cF}{\mathcal{F}}
\newcommand{\cG}{\mathcal{G}}
\newcommand{\cH}{\mathcal{H}}
\newcommand{\cI}{\mathcal{I}}
\newcommand{\cJ}{\mathcal{J}}
\newcommand{\cK}{\mathcal{K}}
\newcommand{\cL}{\mathcal{L}}
\newcommand{\cM}{\mathcal{M}}
\newcommand{\cN}{\mathcal{N}}
\newcommand{\cO}{\mathcal{O}}
\newcommand{\cP}{\mathcal{P}}
\newcommand{\cQ}{\mathcal{Q}}
\newcommand{\cR}{\mathcal{R}}
\newcommand{\cS}{\mathcal{S}}
\newcommand{\cT}{\mathcal{T}}
\newcommand{\cU}{\mathcal{U}}
\newcommand{\cV}{\mathcal{V}}
\newcommand{\cW}{\mathcal{W}}
\newcommand{\cX}{\mathcal{X}}
\newcommand{\cY}{\mathcal{Y}}
\newcommand{\cZ}{\mathcal{Z}}




% --- Common sets ---
\newcommand{\BA}{\mathbb{A}} % predefined for Ångström symbol 
\newcommand{\BB}{\mathbb{B}}
\newcommand{\CC}{\mathbb{C}}
\newcommand{\DD}{\mathbb{D}}
\newcommand{\EE}{\mathbb{E}}
\newcommand{\FF}{\mathbb{F}}
\newcommand{\GG}{\mathbb{G}}
\newcommand{\HH}{\mathbb{H}}
\newcommand{\II}{\mathbb{I}}
\newcommand{\JJ}{\mathbb{J}}
\newcommand{\KK}{\mathbb{K}}
\newcommand{\LL}{\mathbb{L}}
\newcommand{\MM}{\mathbb{M}}
\newcommand{\NN}{\mathbb{N}}
\newcommand{\OO}{\mathbb{O}}
\newcommand{\PP}{\mathbb{P}}
\newcommand{\QQ}{\mathbb{Q}}
\newcommand{\RR}{\mathbb{R}}
\newcommand{\BS}{\mathbb{S}} % predefined for German esszet
\newcommand{\TT}{\mathbb{T}}
\newcommand{\UU}{\mathbb{U}}
\newcommand{\VV}{\mathbb{V}}
\newcommand{\WW}{\mathbb{W}}
\newcommand{\XX}{\mathbb{X}}
\newcommand{\YY}{\mathbb{Y}}
\newcommand{\ZZ}{\mathbb{Z}}

\newcommand{\fa}{\mathfrak{a}}
\newcommand{\fb}{\mathfrak{b}}
\newcommand{\fc}{\mathfrak{c}}
\newcommand{\fd}{\mathfrak{d}}
\newcommand{\fe}{\mathfrak{e}}
\newcommand{\ff}{\mathfrak{f}}
\newcommand{\fg}{\mathfrak{g}}
\newcommand{\fh}{\mathfrak{h}}
\newcommand{\fii}{\mathfrak{i}} % Replaced \fi to avoid breaking TeX
\newcommand{\fj}{\mathfrak{j}}
\newcommand{\fk}{\mathfrak{k}}
\newcommand{\fl}{\mathfrak{l}}
\newcommand{\fm}{\mathfrak{m}}
\newcommand{\fn}{\mathfrak{n}}
\newcommand{\fo}{\mathfrak{o}}
\newcommand{\fp}{\mathfrak{p}}
\newcommand{\fq}{\mathfrak{q}}
\newcommand{\fr}{\mathfrak{r}}
\newcommand{\fs}{\mathfrak{s}}
\newcommand{\ft}{\mathfrak{t}}
\newcommand{\fu}{\mathfrak{u}}
\newcommand{\fv}{\mathfrak{v}}
\newcommand{\fw}{\mathfrak{w}}
\newcommand{\fx}{\mathfrak{x}}
\newcommand{\fy}{\mathfrak{y}}
\newcommand{\fz}{\mathfrak{z}}



\newcommand{\fA}{\mathfrak{A}}
\newcommand{\fB}{\mathfrak{B}}
\newcommand{\fC}{\mathfrak{C}}
\newcommand{\fD}{\mathfrak{D}}
\newcommand{\fE}{\mathfrak{E}}
\newcommand{\fF}{\mathfrak{F}}
\newcommand{\fG}{\mathfrak{G}}
\newcommand{\fH}{\mathfrak{H}}
\newcommand{\fI}{\mathfrak{I}}
\newcommand{\fJ}{\mathfrak{J}}
\newcommand{\fK}{\mathfrak{K}}
\newcommand{\fL}{\mathfrak{L}}
\newcommand{\fM}{\mathfrak{M}}
\newcommand{\fN}{\mathfrak{N}}
\newcommand{\fO}{\mathfrak{O}}
\newcommand{\fP}{\mathfrak{P}}
\newcommand{\fQ}{\mathfrak{Q}}
\newcommand{\fR}{\mathfrak{R}}
\newcommand{\fS}{\mathfrak{S}}
\newcommand{\fT}{\mathfrak{T}}
\newcommand{\fU}{\mathfrak{U}}
\newcommand{\fV}{\mathfrak{V}}
\newcommand{\fW}{\mathfrak{W}}
\newcommand{\fX}{\mathfrak{X}}
\newcommand{\fY}{\mathfrak{Y}}
\newcommand{\fZ}{\mathfrak{Z}}


% --- Math operators ---
\newcommand{\argmax}[1]{\mathop{\mathrm{arg\,max}}_{#1}\,}
\newcommand{\argmin}[1]{\mathop{\mathrm{arg\,min}}_{#1}\,}
\newcommand{\probd}[1]{\textbf{P}(#1)}

% --- Set and tuple utilities ---
\newcommand{\tuple}[1]{\left\langle #1 \right\rangle}
\newcommand{\restr}[2]{#1\!\upharpoonright_{#2}} 
\newcommand{\sint}[1]{\llbracket #1 \rrbracket}  
\newcommand{\setsize}[1]{\Gamma (#1)}            
\newcommand{\inverse}[1]{\ensuremath{{#1}^{-1}}} 
\newcommand{\powerset}[1]{\ensuremath{\mathcal{P}(#1)}} 
\newcommand{\setbuilder}[2]{\ensuremath{\{\,#1 \mid #2\,\}}}

% --- LaTeX3 function syntax macros ---
\ExplSyntaxOn
\NewDocumentCommand{\union}{m}{\ensuremath{\clist_use:nn {#1} { \mathbin{\cup} }}}
\NewDocumentCommand{\intersection}{m}{\ensuremath{\clist_use:nn {#1} { \mathbin{\cap} }}}
\NewDocumentCommand{\compose}{m}{\ensuremath{\clist_use:nn {#1} { \mathbin{\circ} }}}
\NewDocumentCommand{\cartprod}{m}{\ensuremath{\clist_use:nn {#1} { \mathbin{\times} }}}
\ExplSyntaxOff

% --- Total and Partial Functions ---
\newcommand{\func}[3]{\ensuremath{#1 : #2 \to #3}}
\newcommand{\pfunc}[3]{\ensuremath{#1 : #2 \rightharpoonup #3}}

% --- Indexed set operations ---
\newcommand{\bigunion}[1]{\ensuremath{\bigcup_{i \in I} {#1}_{i}}}
\newcommand{\bigintersection}[1]{\ensuremath{\bigcap_{i \in I} {#1}_{i}}}

% --- Common math helpers ---
\newcommand{\identity}[1]{\ensuremath{\operatorname{Id}_{#1}}}
\newcommand{\bigo}[1]{\mathcal{O}(#1)}
\newcommand{\diff}{\mathop{}\!\mathrm{d}}

% ===================================================================
% TIKZ AND DIAGRAMS
% ===================================================================
\usepackage{tikz}
\usetikzlibrary{arrows.meta, positioning, trees, shapes, matrix, calc, decorations.pathreplacing}
\tikzset{
  vertex/.style={circle, draw, minimum size=18pt, inner sep=0pt},
  edge/.style={-},
  arc/.style={->, >=Stealth}
}

% ===================================================================
% CODE, ALGORITHMS, AND BOXES
% ===================================================================
\usepackage{tcolorbox}
\tcbuselibrary{listingsutf8, skins, fitting, theorems, raster, magazine, breakable, xparse}
\usepackage{lipsum,lmodern}


\usepackage{algorithm, algorithmicx, algpseudocode} % Algorithm support

\usepackage[utf8]{inputenc}
\usepackage{textcomp}

% --- Note boxes ---
\newcommand{\commandnote}[1]{%
  \begin{tcolorbox}[standard jigsaw, title=Note]
    #1
  \end{tcolorbox}
}

% --- Highlighted inline text ---
\newcommand{\hltext}[1]{\sethlcolor{gray!10}\hl{#1}}

% --- Code block environment ---
\newtcblisting{codeblock}{
  colback=gray!10,
  colframe=black!40,
  listing only,
  boxrule=0.1pt,
  arc=2pt,
  left=0pt, right=0pt, top=0pt, bottom=0pt,
  listing options={
    breaklines=true,
    basicstyle=\ttfamily,
    literate={\$}{{\textdollar}}1
  }
}

% --- Code block with title ---
\newtcblisting{codeblocktitle}[1][]{
  colback=gray!10,
  colframe=black!40,
  colbacktitle=gray!30,
  coltitle=black,
  listing only,
  boxrule=0.1pt,
  arc=2pt,
  left=0pt, right=0pt, top=0pt, bottom=0pt,
  #1,
  listing options={
    breaklines=true,
    basicstyle=\ttfamily,
    literate={\$}{{\textdollar}}1
  }
}

% --- Captioned code block environment ---
\lstnewenvironment{codeblocklst}[1][]
{
  \lstset{
    backgroundcolor=\color{gray!10},
    frame=single,
    rulecolor=\color{black!10},
    framerule=0.1pt,
    breaklines=true,
    basicstyle=\ttfamily,
    literate={\$}{{\textdollar}}1,
    captionpos=b,
    #1
  }
}{}

% --- Inline code snippet ---
\newcommand{\inlinecode}[1]{%
  \tcbox[on line, fontupper=\ttfamily, size=small, boxrule=0.1pt, colframe=gray!50!black, colback=gray!10!white]{\detokenize{#1}}%
}


% ===================================================================
% BIBLIOGRAPHY AND HYPERLINKS
% ===================================================================
\usepackage[
backend=biber,
style=alphabetic,
sorting=ynt
]{biblatex}
\addbibresource{refs.bib}

\usepackage{hyperref}
\hypersetup{
  colorlinks,
  citecolor=teal,
  linkcolor=blue,
  urlcolor=magenta
}
\usepackage[nameinlink,capitalize]{cleveref}
% ===================================================================
% CUSTOM STRUCTURES AND LABELS
% ===================================================================
% --- Section titles ---
\newcommand{\labeledsection}[2]{\section{#1}\label{#2}}
\newcommand{\minititle}[1]{\subsubsection*{#1}}

% --- Anchors and references ---
\newcommand{\refterm}[2]{\phantomsection \label{#2}{\textbf{#1}}}
\newcommand{\linkterm}[2]{\hyperref[#2]{#1}}

% ===================================================================
% DEFINITIONS AND THEOREMS
% ===================================================================
% --- Definition emphasizer ---
\newcommand{\defemph}[1]{\textbf{#1}}

% --- Section-scoped definitions ---
\newcounter{definition}[section]
\renewcommand{\thedefinition}{\thesection.\arabic{definition}}
\newcommand{\definition}[3]{
  \refstepcounter{definition}
  \phantomsection
  \par\noindent\textbf{Def~\thedefinition.} \textbf{#1: }#2
  \label{#3}\par
}

% --- Section-scoped definition (Numbered, No defineiendum) ---
\newcommand{\anondefinition}[2]{
  \refstepcounter{definition}
  \phantomsection
  \par\noindent\textbf{Def~\thedefinition.} #1
  \label{#2}\par
}

% --- Subsection-scoped definitions ---
\newcounter{subdefinition}[subsection]
\renewcommand{\thesubdefinition}{\thesubsection.\arabic{subdefinition}}
\newcommand{\subdefinition}[3]{
  \refstepcounter{subdefinition}
  \phantomsection
  \par\noindent\textbf{Def~\thesubdefinition.} \textbf{#1: }#2
  \label{#3}\par
}

% --- Subsection-scoped definition (Numbered, No defineiendum) ---
\newcommand{\anonsubdefinition}[2]{
  \refstepcounter{subdefinition}
  \phantomsection
  \par\noindent\textbf{Def~\thesubdefinition.} #1
  \label{#2}\par
}

% --- Shared theorem counter per section ---
\newcounter{statement}[section]
\renewcommand{\thestatement}{\thesection.\arabic{statement}}

% --- Linked theorem-type counters ---
\newcounter{theorem}
\renewcommand{\thetheorem}{\thestatement}
\newcounter{lemma}
\renewcommand{\thelemma}{\thestatement}
\newcounter{corollary}
\renewcommand{\thecorollary}{\thestatement}

% --- Independent example counter ---
\newcounter{example}[section]
\renewcommand{\theexample}{\thesection.\arabic{example}}

% --- Statement macros ---
\newcommand{\Theorem}[2]{
  \stepcounter{statement}\refstepcounter{theorem}
  \phantomsection
  \par\noindent\textbf{Theorem~\thetheorem.} #1
  \label{#2}\par
}
\newcommand{\Lemma}[2]{
  \stepcounter{statement}\refstepcounter{lemma}
  \phantomsection
  \par\noindent\textbf{Lemma~\thelemma.} #1
  \label{#2}\par
}
\newcommand{\Corollary}[2]{
  \stepcounter{statement}\refstepcounter{corollary}
  \phantomsection
  \par\noindent\textbf{Corollary~\thecorollary.} #1
  \label{#2}\par
}
\newcommand{\Example}[2]{
  \refstepcounter{example}
  \phantomsection
  \par\noindent\textbf{Example~\theexample.} #1
  \label{#2}\par
}
\newcommand{\Proof}[1]{
  \par\noindent\textit{Proof.} #1\hfill$\square$\par
}

% ===================================================================
% CLEVEREF CONFIGURATION
% ===================================================================
\crefname{section}{section}{sections}
\Crefname{section}{Section}{Sections}
\crefname{subsection}{section}{sections}
\Crefname{subsection}{Section}{Sections}
\crefname{figure}{figure}{figures}
\Crefname{figure}{Figure}{Figures}
\crefname{table}{table}{tables}
\Crefname{table}{Table}{Tables}
\crefname{equation}{equation}{equations}
\Crefname{equation}{Equation}{Equations}
\crefname{algorithm}{algorithm}{algorithms}
\Crefname{algorithm}{Algorithm}{Algorithms}
\crefname{theorem}{theorem}{theorems}
\Crefname{theorem}{Theorem}{Theorems}
\crefname{lemma}{lemma}{lemmas}
\Crefname{lemma}{Lemma}{Lemmas}
\crefname{corollary}{corollary}{corollaries}
\Crefname{corollary}{Corollary}{Corollaries}
\crefname{example}{example}{examples}
\Crefname{example}{Example}{Examples}
\crefname{codeblock}{listing}{listings}
\Crefname{codeblock}{Listing}{Listings}
\crefname{appendix}{appendix}{appendices}
\Crefname{appendix}{Appendix}{Appendices}
\Crefname{definition}{Def}{Defs}
\Crefname{subdefinition}{Def}{Defs}

% --- Reference shorthands ---
\newcommand{\secref}[1]{\cref{#1}}   \newcommand{\Secref}[1]{\Cref{#1}}
\newcommand{\figref}[1]{\cref{#1}}   \newcommand{\Figref}[1]{\Cref{#1}}
\newcommand{\tabref}[1]{\cref{#1}}   \newcommand{\Tabref}[1]{\Cref{#1}}
\newcommand{\eqrefc}[1]{\cref{#1}}   \newcommand{\Eqref}[1]{\Cref{#1}}
\newcommand{\defref}[1]{\cref{#1}}  \newcommand{\Defref}[1]{\Cref{#1}} 
\newcommand{\subdefref}[1]{\cref{#1}} \newcommand{\Subdefref}[1]{\Cref{#1}}
\newcommand{\thmref}[1]{\cref{#1}}   \newcommand{\Thmref}[1]{\Cref{#1}}
\newcommand{\lemref}[1]{\cref{#1}}   \newcommand{\Lemref}[1]{\Cref{#1}}
\newcommand{\corref}[1]{\cref{#1}}   \newcommand{\Corref}[1]{\Cref{#1}}
\newcommand{\exampleref}[1]{\cref{#1}} \newcommand{\Exampleref}[1]{\Cref{#1}}
\newcommand{\lstref}[1]{\cref{#1}}   \newcommand{\Lstref}[1]{\Cref{#1}}  
\newcommand{\appref}[1]{\cref{#1}}   \newcommand{\Appref}[1]{\Cref{#1}}
\newcommand{\Algref}[1]{\Cref{#1}}


% ------------------------ %
\newcommand{\questionbox}[1]{
\begin{tcolorbox}[
    enhanced,
    attach boxed title to top left={xshift=5mm, yshift=-3mm, yshifttext=-1mm},
    colback=blue!5!white,
    colframe=blue!75!black,
    colbacktitle=red!80!black,
    title=Question,
    fonttitle=\bfseries,
    boxed title style={size=small, colframe=red!50!black} 
]
#1
\end{tcolorbox}
}

\newcommand{\answerbox}[1]{
\begin{tcolorbox}[
    enhanced,
    attach boxed title to top left={xshift=5mm, yshift=-3mm, yshifttext=-1mm},
    colback=green!5!white,      
    colframe=green!40!black,    
    colbacktitle=red!80!black,  
    title=Answer,
    fonttitle=\bfseries,
    boxed title style={size=small, colframe=red!50!black}
]
#1
\end{tcolorbox}
}


\newcommand{\titledbox}[2]{
\begin{tcolorbox}[
    colback=blue!3!white,
    colframe=blue!70!white,
    title=#1]
#2
\end{tcolorbox}
}

\newcommand{\redbox}[1]{
\begin{tcolorbox}[colback=red!5!white,colframe=red!75!black]
#1
\end{tcolorbox}
}

\newcommand{\centeredtitledbox}[2]{
\begin{tcolorbox}[
    enhanced,
    attach boxed title to top center={yshift=-3mm,yshifttext=-1mm},
    colback=blue!5!white,
    colframe=blue!75!black,
    colbacktitle=red!80!black,
    title=#1,
    fonttitle=\bfseries,
    boxed title style={size=small,colframe=red!50!black}
]
#2
\end{tcolorbox}
}

\newcommand{\attentionbox}[1]{
\begin{tcolorbox}[colback=red!5!white,colframe=red!75!black,title=\textbf{ATTENTION}]
#1
\end{tcolorbox}
}

\newcommand{\observationbox}[1]{
\begin{tcolorbox}[
    enhanced,
    attach boxed title to top left={xshift=5mm, yshift=-3mm, yshifttext=-1mm},
    colback=violet!5!white,
    colframe=violet!60!black,
    colbacktitle=purple!80!black,
    title=Observation,
    fonttitle=\bfseries,
    boxed title style={size=small, colframe=purple!50!black}
]
#1
\end{tcolorbox}
}

\newcommand{\ideabox}[1]{
\begin{tcolorbox}[
    enhanced,
    attach boxed title to top left={xshift=5mm, yshift=-3mm, yshifttext=-1mm},
    colback=yellow!8!white,
    colframe=orange!60!black,
    colbacktitle=orange!80!black,
    title=Idea,
    fonttitle=\bfseries,
    boxed title style={size=small, colframe=orange!50!black}
]
#1
\end{tcolorbox}
}


\newcommand{\notationbox}[1]{
\begin{tcolorbox}[
    enhanced,
    attach boxed title to top left={xshift=5mm, yshift=-3mm, yshifttext=-1mm},
    colback=teal!5!white,
    colframe=teal!60!black,
    colbacktitle=teal!80!black,
    title=Notation,
    fonttitle=\bfseries,
    boxed title style={size=small, colframe=teal!50!black}
]
#1
\end{tcolorbox}
}


\usepackage{subcaption}